%% ECON 282E -- Session 4, Deck B5: Worked Example -- the Stochastic RBC Model, Properly
%% The Session 4 analogue of S03_A4.  3.A4 solved the model with every choice made
%% the crudest way; this deck solves the same model with the Session 4 toolkit --
%% a Rouwenhorst chain, interpolation, golden-section maximization and Howard --
%% line by line, with the stage-by-stage measurement.
%% Numbers from tools/figures/s04_numbers.json, key "worked_stochastic".

%% ECON 282E -- Foundations of Macroeconomics (UC Riverside, Fall 2026)
%% Shared Beamer preamble for every deck in slides/S01 ... slides/S10.
%%
%% Usage (a deck lives in slides/S0X/ and is compiled from that folder):
%%   %% ECON 282E -- Foundations of Macroeconomics (UC Riverside, Fall 2026)
%% Shared Beamer preamble for every deck in slides/S01 ... slides/S10.
%%
%% Usage (a deck lives in slides/S0X/ and is compiled from that folder):
%%   %% ECON 282E -- Foundations of Macroeconomics (UC Riverside, Fall 2026)
%% Shared Beamer preamble for every deck in slides/S01 ... slides/S10.
%%
%% Usage (a deck lives in slides/S0X/ and is compiled from that folder):
%%   \input{../preamble/course_preamble.tex}
%%   \decktitle{1}{A1}{Who I Am \& Why This Course}{September 24, 2026}
%%   \begin{document}
%%   \makedecktitle
%%   ...
%%
%% Adapted from Courses/AI-ECON-2026/source/shared_preamble.tex (Metropolis theme,
%% concept/caution/example/takeaway boxes, \sectiondivider). Changes for this course:
%%   * course title block (\decktitle / \makedecktitle) instead of \makebeamertitle
%%   * \zh{...} typesets Chinese (book titles) under pdflatex via CJKutf8
%%   * researchbox for the "Research in the wild" frames
%%   * section pages off; TikZ libraries and the course map loaded here
%% Build with pdflatex only (two passes); tools/build_slides.py does this for you.

\documentclass[10pt,english,aspectratio=169]{beamer}

\usepackage[T1]{fontenc}
\usepackage{textcomp}
\usepackage[utf8]{inputenc}
\usepackage{babel}
\usepackage{CJKutf8}

\usepackage{amsmath, amssymb, amsthm, graphicx}
\usepackage{booktabs, multirow, tabularx}
\usepackage{tikz}
\usetikzlibrary{positioning,arrows.meta,shapes,calc,fit,backgrounds}
\usepackage{pgfplots}
\pgfplotsset{compat=1.16}
\usepackage[most]{tcolorbox}
\tcbuselibrary{skins, breakable}
\usepackage{listings}
\usepackage{xcolor}

\lstset{
  basicstyle=\ttfamily\small,
  keywordstyle=\color{blue!80!black}\bfseries,
  commentstyle=\color{green!50!black}\itshape,
  stringstyle=\color{orange!80!black},
  backgroundcolor=\color{gray!8},
  frame=single,
  rulecolor=\color{gray!40},
  breaklines=true,
  showstringspaces=false,
  numbers=none,
}

\usetheme{metropolis}
\usefonttheme{professionalfonts}
\metroset{sectionpage=none}

\definecolor{LightGray}{RGB}{230,230,230}
\definecolor{RedTitle}{RGB}{0,0,200}
\definecolor{VeryLightGray}{RGB}{245,245,245}
\definecolor{DarkText}{RGB}{0,0,0}
\definecolor{UCRBlue}{RGB}{0,45,106}
\definecolor{UCRGold}{RGB}{241,171,0}

% Stage colours of the course arc (used by the course map and elsewhere)
\colorlet{StageTrunk}{blue!12}
\colorlet{StageBranch}{cyan!14}
\colorlet{StageMachine}{gray!18}
\colorlet{StageWall}{red!14}
\colorlet{StageTools}{orange!20}
\colorlet{StageData}{green!16}
\colorlet{StageAgents}{violet!14}

\hypersetup{bookmarks=false,breaklinks=true,pdfborder={0 0 0},colorlinks=true,
  linkcolor=DarkText,urlcolor=RedTitle}

\setbeamercolor{background canvas}{bg=white}
\setbeamercolor{title}{fg=RedTitle}
\setbeamercolor{frametitle}{bg=VeryLightGray,fg=DarkText}
\setbeamercolor{normal text}{fg=DarkText}
\setbeamercolor{alerted text}{fg=RedTitle}
\setbeamersize{text margin left=5mm, text margin right=5mm}
\addtobeamertemplate{frametitle}{}{\vspace{-0.5em}}

\setbeamertemplate{navigation symbols}{}
\setbeamertemplate{footline}{%
  \hspace{5mm}{\usebeamerfont{page number in head/foot}\color{black!45}%
  ECON 282E \,$\cdot$\, \insertshortdeck}\hfill%
  \usebeamercolor[fg]{page number in head/foot}%
  \usebeamerfont{page number in head/foot}%
  \insertframenumber\,/\,\inserttotalframenumber\kern1em\vskip2pt%
}

% ---------------------------------------------------------------------------
% Chinese text under pdflatex (book titles etc.)
\newcommand{\zh}[1]{\begin{CJK*}{UTF8}{gbsn}#1\end{CJK*}}

% ---------------------------------------------------------------------------
% Deck title block
%   \decktitle{session}{deck id}{title}{date}
\newcommand{\insertshortdeck}{}
\newcommand{\decksession}{}
\newcommand{\deckid}{}
\newcommand{\decktitletext}{}
\newcommand{\deckdate}{}
\newcommand{\decktitle}[4]{%
  \renewcommand{\decksession}{#1}%
  \renewcommand{\deckid}{#2}%
  \renewcommand{\decktitletext}{#3}%
  \renewcommand{\deckdate}{#4}%
  \renewcommand{\insertshortdeck}{Session #1 \,$\cdot$\, Deck #1.#2}%
  \title{#3}%
  \author{Zhigang Feng}%
  \date{#4}%
}
\newcommand{\makedecktitle}{%
  \begin{frame}[plain,noframenumbering]
    \begin{tikzpicture}[remember picture,overlay]
      \fill[UCRBlue] (current page.north west) rectangle ([yshift=-0.55cm]current page.north east);
      \fill[UCRGold] ([yshift=-0.55cm]current page.north west) rectangle ([yshift=-0.62cm]current page.north east);
      \node[anchor=west, text=white, font=\small\bfseries] at ([xshift=5mm,yshift=-0.275cm]current page.north west)
        {ECON 282E \,$\cdot$\, Foundations of Macroeconomics};
      \node[anchor=east, text=white, font=\small] at ([xshift=-5mm,yshift=-0.275cm]current page.north east)
        {UC Riverside \,$\cdot$\, Fall 2026};
    \end{tikzpicture}
    \vspace*{1.1cm}
    {\usebeamerfont{title}\color{black!55}\large Session \decksession\ \,$\cdot$\, Deck \decksession.\deckid\par}
    \vspace{0.25cm}
    {\usebeamerfont{title}\color{RedTitle}\LARGE\bfseries \decktitletext\par}
    \vspace{0.7cm}
    {\large Zhigang Feng\par}
    \vspace{0.15cm}
    {\color{black!65}\deckdate\par}
    \vfill
    {\footnotesize\itshape\color{black!55}Build it on paper $\;\to\;$ solve it on the machine $\;\to\;$ push it to the frontier with AI\par}
  \end{frame}}

% ---------------------------------------------------------------------------
% Section divider (plain frame, not counted)
\newcommand{\sectiondivider}[2]{%
\begin{frame}[plain,noframenumbering]
\begin{center}
\vfill
{\Large\textbf{#1}\par}
\vspace{0.35cm}
{\normalsize #2\par}
\vfill
\end{center}
\end{frame}
}

% ---------------------------------------------------------------------------
% Boxes
\newtcolorbox{conceptbox}[1][]{colback=blue!3, colframe=RedTitle!55, boxrule=0.35mm,
  arc=1mm, left=2mm, right=2mm, top=1mm, bottom=1mm, #1}
\newtcolorbox{cautionbox}[1][]{colback=red!3, colframe=red!50!black, boxrule=0.35mm,
  arc=1mm, left=2mm, right=2mm, top=1mm, bottom=1mm, #1}
\newtcolorbox{examplebox}[1][]{colback=green!3, colframe=green!45!black, boxrule=0.35mm,
  arc=1mm, left=2mm, right=2mm, top=1mm, bottom=1mm, #1}
\newtcolorbox{takeawaybox}[1][]{colback=VeryLightGray, colframe=RedTitle!55, boxrule=0.35mm,
  arc=1mm, left=2mm, right=2mm, top=1mm, bottom=1mm, #1}
\newtcolorbox{researchbox}[1][]{enhanced, colback=violet!4, colframe=violet!60!black,
  boxrule=0.35mm, arc=1mm, left=2mm, right=2mm, top=1mm, bottom=1mm,
  fonttitle=\bfseries\small, title={Research in the wild}, #1}

\newcommand{\compacttables}{\renewcommand{\arraystretch}{1.15}}
\newcommand{\src}[1]{{\tiny\color{black!55}#1}}

% ---------------------------------------------------------------------------
\input{../preamble/course_map.tex}

%%   \decktitle{1}{A1}{Who I Am \& Why This Course}{September 24, 2026}
%%   \begin{document}
%%   \makedecktitle
%%   ...
%%
%% Adapted from Courses/AI-ECON-2026/source/shared_preamble.tex (Metropolis theme,
%% concept/caution/example/takeaway boxes, \sectiondivider). Changes for this course:
%%   * course title block (\decktitle / \makedecktitle) instead of \makebeamertitle
%%   * \zh{...} typesets Chinese (book titles) under pdflatex via CJKutf8
%%   * researchbox for the "Research in the wild" frames
%%   * section pages off; TikZ libraries and the course map loaded here
%% Build with pdflatex only (two passes); tools/build_slides.py does this for you.

\documentclass[10pt,english,aspectratio=169]{beamer}

\usepackage[T1]{fontenc}
\usepackage{textcomp}
\usepackage[utf8]{inputenc}
\usepackage{babel}
\usepackage{CJKutf8}

\usepackage{amsmath, amssymb, amsthm, graphicx}
\usepackage{booktabs, multirow, tabularx}
\usepackage{tikz}
\usetikzlibrary{positioning,arrows.meta,shapes,calc,fit,backgrounds}
\usepackage{pgfplots}
\pgfplotsset{compat=1.16}
\usepackage[most]{tcolorbox}
\tcbuselibrary{skins, breakable}
\usepackage{listings}
\usepackage{xcolor}

\lstset{
  basicstyle=\ttfamily\small,
  keywordstyle=\color{blue!80!black}\bfseries,
  commentstyle=\color{green!50!black}\itshape,
  stringstyle=\color{orange!80!black},
  backgroundcolor=\color{gray!8},
  frame=single,
  rulecolor=\color{gray!40},
  breaklines=true,
  showstringspaces=false,
  numbers=none,
}

\usetheme{metropolis}
\usefonttheme{professionalfonts}
\metroset{sectionpage=none}

\definecolor{LightGray}{RGB}{230,230,230}
\definecolor{RedTitle}{RGB}{0,0,200}
\definecolor{VeryLightGray}{RGB}{245,245,245}
\definecolor{DarkText}{RGB}{0,0,0}
\definecolor{UCRBlue}{RGB}{0,45,106}
\definecolor{UCRGold}{RGB}{241,171,0}

% Stage colours of the course arc (used by the course map and elsewhere)
\colorlet{StageTrunk}{blue!12}
\colorlet{StageBranch}{cyan!14}
\colorlet{StageMachine}{gray!18}
\colorlet{StageWall}{red!14}
\colorlet{StageTools}{orange!20}
\colorlet{StageData}{green!16}
\colorlet{StageAgents}{violet!14}

\hypersetup{bookmarks=false,breaklinks=true,pdfborder={0 0 0},colorlinks=true,
  linkcolor=DarkText,urlcolor=RedTitle}

\setbeamercolor{background canvas}{bg=white}
\setbeamercolor{title}{fg=RedTitle}
\setbeamercolor{frametitle}{bg=VeryLightGray,fg=DarkText}
\setbeamercolor{normal text}{fg=DarkText}
\setbeamercolor{alerted text}{fg=RedTitle}
\setbeamersize{text margin left=5mm, text margin right=5mm}
\addtobeamertemplate{frametitle}{}{\vspace{-0.5em}}

\setbeamertemplate{navigation symbols}{}
\setbeamertemplate{footline}{%
  \hspace{5mm}{\usebeamerfont{page number in head/foot}\color{black!45}%
  ECON 282E \,$\cdot$\, \insertshortdeck}\hfill%
  \usebeamercolor[fg]{page number in head/foot}%
  \usebeamerfont{page number in head/foot}%
  \insertframenumber\,/\,\inserttotalframenumber\kern1em\vskip2pt%
}

% ---------------------------------------------------------------------------
% Chinese text under pdflatex (book titles etc.)
\newcommand{\zh}[1]{\begin{CJK*}{UTF8}{gbsn}#1\end{CJK*}}

% ---------------------------------------------------------------------------
% Deck title block
%   \decktitle{session}{deck id}{title}{date}
\newcommand{\insertshortdeck}{}
\newcommand{\decksession}{}
\newcommand{\deckid}{}
\newcommand{\decktitletext}{}
\newcommand{\deckdate}{}
\newcommand{\decktitle}[4]{%
  \renewcommand{\decksession}{#1}%
  \renewcommand{\deckid}{#2}%
  \renewcommand{\decktitletext}{#3}%
  \renewcommand{\deckdate}{#4}%
  \renewcommand{\insertshortdeck}{Session #1 \,$\cdot$\, Deck #1.#2}%
  \title{#3}%
  \author{Zhigang Feng}%
  \date{#4}%
}
\newcommand{\makedecktitle}{%
  \begin{frame}[plain,noframenumbering]
    \begin{tikzpicture}[remember picture,overlay]
      \fill[UCRBlue] (current page.north west) rectangle ([yshift=-0.55cm]current page.north east);
      \fill[UCRGold] ([yshift=-0.55cm]current page.north west) rectangle ([yshift=-0.62cm]current page.north east);
      \node[anchor=west, text=white, font=\small\bfseries] at ([xshift=5mm,yshift=-0.275cm]current page.north west)
        {ECON 282E \,$\cdot$\, Foundations of Macroeconomics};
      \node[anchor=east, text=white, font=\small] at ([xshift=-5mm,yshift=-0.275cm]current page.north east)
        {UC Riverside \,$\cdot$\, Fall 2026};
    \end{tikzpicture}
    \vspace*{1.1cm}
    {\usebeamerfont{title}\color{black!55}\large Session \decksession\ \,$\cdot$\, Deck \decksession.\deckid\par}
    \vspace{0.25cm}
    {\usebeamerfont{title}\color{RedTitle}\LARGE\bfseries \decktitletext\par}
    \vspace{0.7cm}
    {\large Zhigang Feng\par}
    \vspace{0.15cm}
    {\color{black!65}\deckdate\par}
    \vfill
    {\footnotesize\itshape\color{black!55}Build it on paper $\;\to\;$ solve it on the machine $\;\to\;$ push it to the frontier with AI\par}
  \end{frame}}

% ---------------------------------------------------------------------------
% Section divider (plain frame, not counted)
\newcommand{\sectiondivider}[2]{%
\begin{frame}[plain,noframenumbering]
\begin{center}
\vfill
{\Large\textbf{#1}\par}
\vspace{0.35cm}
{\normalsize #2\par}
\vfill
\end{center}
\end{frame}
}

% ---------------------------------------------------------------------------
% Boxes
\newtcolorbox{conceptbox}[1][]{colback=blue!3, colframe=RedTitle!55, boxrule=0.35mm,
  arc=1mm, left=2mm, right=2mm, top=1mm, bottom=1mm, #1}
\newtcolorbox{cautionbox}[1][]{colback=red!3, colframe=red!50!black, boxrule=0.35mm,
  arc=1mm, left=2mm, right=2mm, top=1mm, bottom=1mm, #1}
\newtcolorbox{examplebox}[1][]{colback=green!3, colframe=green!45!black, boxrule=0.35mm,
  arc=1mm, left=2mm, right=2mm, top=1mm, bottom=1mm, #1}
\newtcolorbox{takeawaybox}[1][]{colback=VeryLightGray, colframe=RedTitle!55, boxrule=0.35mm,
  arc=1mm, left=2mm, right=2mm, top=1mm, bottom=1mm, #1}
\newtcolorbox{researchbox}[1][]{enhanced, colback=violet!4, colframe=violet!60!black,
  boxrule=0.35mm, arc=1mm, left=2mm, right=2mm, top=1mm, bottom=1mm,
  fonttitle=\bfseries\small, title={Research in the wild}, #1}

\newcommand{\compacttables}{\renewcommand{\arraystretch}{1.15}}
\newcommand{\src}[1]{{\tiny\color{black!55}#1}}

% ---------------------------------------------------------------------------
%% Course map: the recurring "you are here" slide for ECON 282E.
%% Loaded by course_preamble.tex. Pattern follows AI-ECON-2026 lec01_map.tex, but the map
%% shows the whole ten-session course rather than one lecture.
%%
%%   \CourseMap{s}{label}   s = 1..10 highlights session s and prints "you are here: label"
%%                          (e.g. \CourseMap{1}{1.A2}); earlier sessions get a check mark.
%%   \CourseMap{0}{}        full overview, nothing highlighted.
%%
%% Note: TikZ option lists are comma-split before TeX conditionals run, so every \ifnum
%% below sits outside [...] option lists.

\newcommand{\CourseMap}[2]{%
\begin{center}
\begin{tikzpicture}[
    sess/.style={rectangle, rounded corners=2pt, draw=black!45, minimum width=1.34cm,
        minimum height=1.12cm, text width=1.26cm, align=center, inner sep=1pt},
    stage/.style={font=\tiny\bfseries, text=black!70},
]
  \foreach \i/\d/\t/\c in {%
      1/Sep 24/Intro \\ Growth/StageTrunk,
      2/Oct 1/HA $\cdot$ OLG \\ Policy/StageBranch,
      3/Oct 8/Num.\ DP \\ Python/StageMachine,
      4/Oct 15/Numerics \\ I/StageMachine,
      5/Oct 22/Numerics \\ II/StageMachine,
      6/Oct 29/The wall \\ \textbf{Midterm}/StageWall,
      7/Nov 5/ML \\ Deep nets/StageTools,
      8/Nov 12/RL \\ HA $+$ DL/StageTools,
      9/Nov 19/Text \\ LLMs/StageData,
      10/Dec 3/Agents \\ \textbf{Projects}/StageAgents} {
    \node[sess, fill=\c] (n\i) at ({1.46*(\i-1)}, 0)
      {{\scriptsize\bfseries S\i}\\[-1pt]{\tiny\color{black!60}\d}\\[1pt]{\tiny \t}};
    \ifnum\i<#1
      \node[font=\scriptsize, text=green!45!black] at ([xshift=-2.5mm,yshift=-2.2mm]n\i.north east) {$\checkmark$};
    \fi
    \ifnum\i=#1
      \draw[RedTitle, very thick, rounded corners=2pt]
        ([xshift=-0.6mm,yshift=-0.6mm]n\i.south west) rectangle ([xshift=0.6mm,yshift=0.6mm]n\i.north east);
      % keep the label inside the picture at both ends of the row
      \ifnum\i<3
        \node[font=\scriptsize\bfseries, text=RedTitle, anchor=south west] at ([yshift=1.2mm]n\i.north west)
          {$\blacktriangledown$\,you are here: #2};
      \else\ifnum\i>8
        \node[font=\scriptsize\bfseries, text=RedTitle, anchor=south east] at ([yshift=1.2mm]n\i.north east)
          {you are here: #2\,$\blacktriangledown$};
      \else
        \node[font=\scriptsize\bfseries, text=RedTitle, anchor=south] at ([yshift=1.2mm]n\i.north)
          {you are here: #2\,$\blacktriangledown$};
      \fi\fi
    \fi
  }
  % stage band under the sessions
  \foreach \a/\b/\lab/\c in {1/1/trunk/StageTrunk, 2/2/branches/StageBranch,
      3/5/the machine $\cdot$ classical methods/StageMachine, 6/6/the wall/StageWall,
      7/8/new tools/StageTools, 9/9/new data/StageData, 10/10/colleagues/StageAgents} {
    \fill[\c] ([yshift=-1.5mm]n\a.south west) rectangle ([yshift=-4.6mm]n\b.south east);
    \path ([yshift=-1.5mm]n\a.south west) -- ([yshift=-4.6mm]n\b.south east)
      node[midway, stage] {\lab};
  }
  \node[font=\footnotesize\itshape, text=black!70] at ([yshift=-9.5mm]$(n1.south)!0.5!(n10.south)$)
    {Build it on paper $\;\to\;$ solve it on the machine $\;\to\;$ push it to the frontier with AI};
\end{tikzpicture}
\end{center}
}


%%   \decktitle{1}{A1}{Who I Am \& Why This Course}{September 24, 2026}
%%   \begin{document}
%%   \makedecktitle
%%   ...
%%
%% Adapted from Courses/AI-ECON-2026/source/shared_preamble.tex (Metropolis theme,
%% concept/caution/example/takeaway boxes, \sectiondivider). Changes for this course:
%%   * course title block (\decktitle / \makedecktitle) instead of \makebeamertitle
%%   * \zh{...} typesets Chinese (book titles) under pdflatex via CJKutf8
%%   * researchbox for the "Research in the wild" frames
%%   * section pages off; TikZ libraries and the course map loaded here
%% Build with pdflatex only (two passes); tools/build_slides.py does this for you.

\documentclass[10pt,english,aspectratio=169]{beamer}

\usepackage[T1]{fontenc}
\usepackage{textcomp}
\usepackage[utf8]{inputenc}
\usepackage{babel}
\usepackage{CJKutf8}

\usepackage{amsmath, amssymb, amsthm, graphicx}
\usepackage{booktabs, multirow, tabularx}
\usepackage{tikz}
\usetikzlibrary{positioning,arrows.meta,shapes,calc,fit,backgrounds}
\usepackage{pgfplots}
\pgfplotsset{compat=1.16}
\usepackage[most]{tcolorbox}
\tcbuselibrary{skins, breakable}
\usepackage{listings}
\usepackage{xcolor}

\lstset{
  basicstyle=\ttfamily\small,
  keywordstyle=\color{blue!80!black}\bfseries,
  commentstyle=\color{green!50!black}\itshape,
  stringstyle=\color{orange!80!black},
  backgroundcolor=\color{gray!8},
  frame=single,
  rulecolor=\color{gray!40},
  breaklines=true,
  showstringspaces=false,
  numbers=none,
}

\usetheme{metropolis}
\usefonttheme{professionalfonts}
\metroset{sectionpage=none}

\definecolor{LightGray}{RGB}{230,230,230}
\definecolor{RedTitle}{RGB}{0,0,200}
\definecolor{VeryLightGray}{RGB}{245,245,245}
\definecolor{DarkText}{RGB}{0,0,0}
\definecolor{UCRBlue}{RGB}{0,45,106}
\definecolor{UCRGold}{RGB}{241,171,0}

% Stage colours of the course arc (used by the course map and elsewhere)
\colorlet{StageTrunk}{blue!12}
\colorlet{StageBranch}{cyan!14}
\colorlet{StageMachine}{gray!18}
\colorlet{StageWall}{red!14}
\colorlet{StageTools}{orange!20}
\colorlet{StageData}{green!16}
\colorlet{StageAgents}{violet!14}

\hypersetup{bookmarks=false,breaklinks=true,pdfborder={0 0 0},colorlinks=true,
  linkcolor=DarkText,urlcolor=RedTitle}

\setbeamercolor{background canvas}{bg=white}
\setbeamercolor{title}{fg=RedTitle}
\setbeamercolor{frametitle}{bg=VeryLightGray,fg=DarkText}
\setbeamercolor{normal text}{fg=DarkText}
\setbeamercolor{alerted text}{fg=RedTitle}
\setbeamersize{text margin left=5mm, text margin right=5mm}
\addtobeamertemplate{frametitle}{}{\vspace{-0.5em}}

\setbeamertemplate{navigation symbols}{}
\setbeamertemplate{footline}{%
  \hspace{5mm}{\usebeamerfont{page number in head/foot}\color{black!45}%
  ECON 282E \,$\cdot$\, \insertshortdeck}\hfill%
  \usebeamercolor[fg]{page number in head/foot}%
  \usebeamerfont{page number in head/foot}%
  \insertframenumber\,/\,\inserttotalframenumber\kern1em\vskip2pt%
}

% ---------------------------------------------------------------------------
% Chinese text under pdflatex (book titles etc.)
\newcommand{\zh}[1]{\begin{CJK*}{UTF8}{gbsn}#1\end{CJK*}}

% ---------------------------------------------------------------------------
% Deck title block
%   \decktitle{session}{deck id}{title}{date}
\newcommand{\insertshortdeck}{}
\newcommand{\decksession}{}
\newcommand{\deckid}{}
\newcommand{\decktitletext}{}
\newcommand{\deckdate}{}
\newcommand{\decktitle}[4]{%
  \renewcommand{\decksession}{#1}%
  \renewcommand{\deckid}{#2}%
  \renewcommand{\decktitletext}{#3}%
  \renewcommand{\deckdate}{#4}%
  \renewcommand{\insertshortdeck}{Session #1 \,$\cdot$\, Deck #1.#2}%
  \title{#3}%
  \author{Zhigang Feng}%
  \date{#4}%
}
\newcommand{\makedecktitle}{%
  \begin{frame}[plain,noframenumbering]
    \begin{tikzpicture}[remember picture,overlay]
      \fill[UCRBlue] (current page.north west) rectangle ([yshift=-0.55cm]current page.north east);
      \fill[UCRGold] ([yshift=-0.55cm]current page.north west) rectangle ([yshift=-0.62cm]current page.north east);
      \node[anchor=west, text=white, font=\small\bfseries] at ([xshift=5mm,yshift=-0.275cm]current page.north west)
        {ECON 282E \,$\cdot$\, Foundations of Macroeconomics};
      \node[anchor=east, text=white, font=\small] at ([xshift=-5mm,yshift=-0.275cm]current page.north east)
        {UC Riverside \,$\cdot$\, Fall 2026};
    \end{tikzpicture}
    \vspace*{1.1cm}
    {\usebeamerfont{title}\color{black!55}\large Session \decksession\ \,$\cdot$\, Deck \decksession.\deckid\par}
    \vspace{0.25cm}
    {\usebeamerfont{title}\color{RedTitle}\LARGE\bfseries \decktitletext\par}
    \vspace{0.7cm}
    {\large Zhigang Feng\par}
    \vspace{0.15cm}
    {\color{black!65}\deckdate\par}
    \vfill
    {\footnotesize\itshape\color{black!55}Build it on paper $\;\to\;$ solve it on the machine $\;\to\;$ push it to the frontier with AI\par}
  \end{frame}}

% ---------------------------------------------------------------------------
% Section divider (plain frame, not counted)
\newcommand{\sectiondivider}[2]{%
\begin{frame}[plain,noframenumbering]
\begin{center}
\vfill
{\Large\textbf{#1}\par}
\vspace{0.35cm}
{\normalsize #2\par}
\vfill
\end{center}
\end{frame}
}

% ---------------------------------------------------------------------------
% Boxes
\newtcolorbox{conceptbox}[1][]{colback=blue!3, colframe=RedTitle!55, boxrule=0.35mm,
  arc=1mm, left=2mm, right=2mm, top=1mm, bottom=1mm, #1}
\newtcolorbox{cautionbox}[1][]{colback=red!3, colframe=red!50!black, boxrule=0.35mm,
  arc=1mm, left=2mm, right=2mm, top=1mm, bottom=1mm, #1}
\newtcolorbox{examplebox}[1][]{colback=green!3, colframe=green!45!black, boxrule=0.35mm,
  arc=1mm, left=2mm, right=2mm, top=1mm, bottom=1mm, #1}
\newtcolorbox{takeawaybox}[1][]{colback=VeryLightGray, colframe=RedTitle!55, boxrule=0.35mm,
  arc=1mm, left=2mm, right=2mm, top=1mm, bottom=1mm, #1}
\newtcolorbox{researchbox}[1][]{enhanced, colback=violet!4, colframe=violet!60!black,
  boxrule=0.35mm, arc=1mm, left=2mm, right=2mm, top=1mm, bottom=1mm,
  fonttitle=\bfseries\small, title={Research in the wild}, #1}

\newcommand{\compacttables}{\renewcommand{\arraystretch}{1.15}}
\newcommand{\src}[1]{{\tiny\color{black!55}#1}}

% ---------------------------------------------------------------------------
%% Course map: the recurring "you are here" slide for ECON 282E.
%% Loaded by course_preamble.tex. Pattern follows AI-ECON-2026 lec01_map.tex, but the map
%% shows the whole ten-session course rather than one lecture.
%%
%%   \CourseMap{s}{label}   s = 1..10 highlights session s and prints "you are here: label"
%%                          (e.g. \CourseMap{1}{1.A2}); earlier sessions get a check mark.
%%   \CourseMap{0}{}        full overview, nothing highlighted.
%%
%% Note: TikZ option lists are comma-split before TeX conditionals run, so every \ifnum
%% below sits outside [...] option lists.

\newcommand{\CourseMap}[2]{%
\begin{center}
\begin{tikzpicture}[
    sess/.style={rectangle, rounded corners=2pt, draw=black!45, minimum width=1.34cm,
        minimum height=1.12cm, text width=1.26cm, align=center, inner sep=1pt},
    stage/.style={font=\tiny\bfseries, text=black!70},
]
  \foreach \i/\d/\t/\c in {%
      1/Sep 24/Intro \\ Growth/StageTrunk,
      2/Oct 1/HA $\cdot$ OLG \\ Policy/StageBranch,
      3/Oct 8/Num.\ DP \\ Python/StageMachine,
      4/Oct 15/Numerics \\ I/StageMachine,
      5/Oct 22/Numerics \\ II/StageMachine,
      6/Oct 29/The wall \\ \textbf{Midterm}/StageWall,
      7/Nov 5/ML \\ Deep nets/StageTools,
      8/Nov 12/RL \\ HA $+$ DL/StageTools,
      9/Nov 19/Text \\ LLMs/StageData,
      10/Dec 3/Agents \\ \textbf{Projects}/StageAgents} {
    \node[sess, fill=\c] (n\i) at ({1.46*(\i-1)}, 0)
      {{\scriptsize\bfseries S\i}\\[-1pt]{\tiny\color{black!60}\d}\\[1pt]{\tiny \t}};
    \ifnum\i<#1
      \node[font=\scriptsize, text=green!45!black] at ([xshift=-2.5mm,yshift=-2.2mm]n\i.north east) {$\checkmark$};
    \fi
    \ifnum\i=#1
      \draw[RedTitle, very thick, rounded corners=2pt]
        ([xshift=-0.6mm,yshift=-0.6mm]n\i.south west) rectangle ([xshift=0.6mm,yshift=0.6mm]n\i.north east);
      % keep the label inside the picture at both ends of the row
      \ifnum\i<3
        \node[font=\scriptsize\bfseries, text=RedTitle, anchor=south west] at ([yshift=1.2mm]n\i.north west)
          {$\blacktriangledown$\,you are here: #2};
      \else\ifnum\i>8
        \node[font=\scriptsize\bfseries, text=RedTitle, anchor=south east] at ([yshift=1.2mm]n\i.north east)
          {you are here: #2\,$\blacktriangledown$};
      \else
        \node[font=\scriptsize\bfseries, text=RedTitle, anchor=south] at ([yshift=1.2mm]n\i.north)
          {you are here: #2\,$\blacktriangledown$};
      \fi\fi
    \fi
  }
  % stage band under the sessions
  \foreach \a/\b/\lab/\c in {1/1/trunk/StageTrunk, 2/2/branches/StageBranch,
      3/5/the machine $\cdot$ classical methods/StageMachine, 6/6/the wall/StageWall,
      7/8/new tools/StageTools, 9/9/new data/StageData, 10/10/colleagues/StageAgents} {
    \fill[\c] ([yshift=-1.5mm]n\a.south west) rectangle ([yshift=-4.6mm]n\b.south east);
    \path ([yshift=-1.5mm]n\a.south west) -- ([yshift=-4.6mm]n\b.south east)
      node[midway, stage] {\lab};
  }
  \node[font=\footnotesize\itshape, text=black!70] at ([yshift=-9.5mm]$(n1.south)!0.5!(n10.south)$)
    {Build it on paper $\;\to\;$ solve it on the machine $\;\to\;$ push it to the frontier with AI};
\end{tikzpicture}
\end{center}
}


\graphicspath{{./figures/}}

\lstset{language=Python, basicstyle=\ttfamily\scriptsize, frame=none,
        backgroundcolor=\color{gray!8}, aboveskip=2pt, belowskip=2pt}

\decktitle{4}{B5}{Worked Example: The Stochastic RBC Model, Properly}{Thursday, October 15, 2026}

\begin{document}

\makedecktitle

%=============================================================================
\begin{frame}{Where We Are}
\CourseMap{4}{4.B5}
\vspace{-1mm}
\begin{center}
\small \textbf{Block B:} \textbf{4.B1} approximation and interpolation $\;\cdot\;$ \textbf{4.B2} quadrature $\;\cdot\;$ \textbf{4.B3} discretizing an AR(1) $\;\cdot\;$ \textbf{4.B4} fast and accurate $\;\cdot\;$ \textbf{4.B5} the worked example, end to end.
\end{center}
\end{frame}

%=============================================================================
\begin{frame}{The Same Deck as 3.A4, With the Tools Turned On}
\begin{columns}[T]
\begin{column}{0.53\textwidth}
\small
In 3.A4 every one of the six decisions was made the crudest way: a hand-written two-state chain, no interpolation, a grid search for the maximum, and a full sweep every iteration. We re-run \emph{that solver} here, at this session's quarterly calibration, to get an honest ``before'' column.
\medskip

Then the same five steps, with this session's replacements in place:
\medskip
\begin{center}\footnotesize
\begin{tabularx}{\linewidth}{@{}>{\raggedright\arraybackslash}p{2.3cm} >{\raggedright\arraybackslash}X@{}}
\toprule
\textbf{3.A4} & \textbf{here} \\
\midrule
$\Pi$ by hand & \textbf{Rouwenhorst} (4.B3) \\
no interpolation & \textbf{linear}, which keeps concavity (4.B1) \\
grid search & \textbf{golden section} (4.A2) \\
full sweeps & \textbf{Howard}, $n_H=20$ (4.B4) \\
\bottomrule
\end{tabularx}
\end{center}
\end{column}
\begin{column}{0.44\textwidth}
\begin{conceptbox}[title=\textbf{The model}]\small
\[
  v(k,z)=\max_{k'}\Big\{\ln c+\beta\,\mathbb{E}\big[v(k',z')\mid z\big]\Big\}
\]
\[
  c+k'=zk^{\alpha}+(1-\delta)k, \quad \ln z'=\rho\ln z+\sigma\varepsilon' .
\]
Quarterly: $\alpha=0.33$, $\beta=0.99$, $\delta=0.025$, $\rho=0.95$, $\sigma=0.007$.
\end{conceptbox}
\medskip
\begin{examplebox}\footnotesize
No closed-form policy. The benchmarks are the steady state, the chain's own moments, and the Euler residual.
\end{examplebox}
\end{column}
\end{columns}
\end{frame}

%=============================================================================
\begin{frame}[fragile]{Step 1: The Chain, and Checking It Before Anything Else}
\begin{columns}[T]
\begin{column}{0.54\textwidth}
\begin{lstlisting}
def rouwenhorst(n, rho, sigma):
    p = (1 + rho) / 2
    P = np.array([[p, 1-p], [1-p, p]])
    for k in range(3, n+1):
        Z = np.zeros((k, k))
        Z[:-1,:-1] += p*P;     Z[:-1,1:] += (1-p)*P
        Z[1:, :-1] += (1-p)*P; Z[1:, 1:] += p*P
        Z[1:-1, :] /= 2.0
        P = Z
    sz  = sigma / np.sqrt(1 - rho**2)
    psi = sz * np.sqrt(n - 1)
    return np.linspace(-psi, psi, n), P

logz, Pi = rouwenhorst(7, 0.95, 0.007)
z = np.exp(logz)
\end{lstlisting}
\end{column}
\begin{column}{0.43\textwidth}
\small
\begin{center}\footnotesize
\begin{tabularx}{\linewidth}{@{}>{\raggedright\arraybackslash}X r@{}}
\toprule
implied $\rho$ & 0.950000 \\
implied $\sigma_z$ & 0.02241794 \\
target $\sigma_z$ & 0.02241794 \\
$z$ range & $[0.947,\,1.056]$ \\
\bottomrule
\end{tabularx}
\end{center}
\vspace{1mm}
\begin{takeawaybox}\footnotesize
\textbf{Check the chain before you solve anything.} Three lines, and it rules out the single most common silent error in this literature --- a shock process that is not the one you wrote down.
\end{takeawaybox}
\end{column}
\end{columns}
\end{frame}

%=============================================================================
\begin{frame}[fragile]{Step 2: The Benchmark, and the Grid}
\begin{columns}[T]
\begin{column}{0.53\textwidth}
\begin{lstlisting}
alpha, beta, delta = 0.33, 0.99, 0.025

# the one thing we know exactly
kss = ((1/beta - 1 + delta)/alpha)**(1/(alpha-1))
# 28.348419061048443

n_k = 300
kg  = np.linspace(0.5*kss, 1.5*kss, n_k)
h   = kg[1] - kg[0]          # 0.0948

# cash on hand, every (state, shock)
y = (z[None,:]*kg[:,None]**alpha
     + (1-delta)*kg[:,None])
\end{lstlisting}
\end{column}
\begin{column}{0.44\textwidth}
\small
$k^{*}=28.3484$, $y^{*}=3.0153$, $c^{*}=2.3066$, and $K/Y=9.40$ at a quarterly frequency --- about $2.35$ annual, which is the number to sanity-check against.
\medskip
\begin{cautionbox}\footnotesize
The grid is $\pm50\%$ around $k^{*}$. Print $\max_i g(k_i)/k_N$ when you are done: if it is 1, the policy is pinned at the top of the grid and the answer is censored, not solved.
\end{cautionbox}
\end{column}
\end{columns}
\end{frame}

%=============================================================================
\begin{frame}[fragile]{Step 3: The Inner Maximization Is No Longer a Search}
\begin{columns}[T]
\begin{column}{0.55\textwidth}
\begin{lstlisting}
PHI = (np.sqrt(5) - 1) / 2

def rhs(kp, yj, EVj):    # Bellman RHS, off grid
    c = np.maximum(yj - kp, 1e-12)   # clamp
    return np.log(c) + beta*np.interp(kp,kg,EVj)

# golden section, vectorised over all k at once
yj, EVj = y[:, j], EV[:, j]
a = np.full(n_k, kg[0])
b = np.minimum(yj - 1e-8, kg[-1])
c_, d_ = b - PHI*(b-a), a + PHI*(b-a)
fc, fd = rhs(c_, yj, EVj), rhs(d_, yj, EVj)
for _ in range(40):
    L  = fc > fd
    b  = np.where(L, d_, b); a = np.where(L, a, c_)
    c_, d_ = b - PHI*(b-a), a + PHI*(b-a)
    fc, fd = rhs(c_, yj, EVj), rhs(d_, yj, EVj)
g[:, j] = 0.5*(a + b)
\end{lstlisting}
\end{column}
\begin{column}{0.42\textwidth}
\small
Three of this session's pieces, in one block:
\medskip
\begin{itemize}\footnotesize
  \item \texttt{np.interp} is \textbf{4.B1's} linear interpolation --- chosen because it keeps $\hat v$ concave;
  \item that concavity is what makes the objective \textbf{unimodal};
  \item which is what licenses \textbf{4.A2's} golden section to shrink the bracket safely.
\end{itemize}
\medskip
\begin{cautionbox}\footnotesize
Swap in a cubic spline and the first bullet fails, so the third one does too. \textbf{The interpolation choice is not cosmetic} --- it is a hypothesis of the maximizer.
\end{cautionbox}
\end{column}
\end{columns}
\end{frame}

%=============================================================================
\begin{frame}[fragile]{Step 4: The Expectation, and Howard Around the Outside}
\begin{columns}[T]
\begin{column}{0.53\textwidth}
\begin{lstlisting}
V = np.zeros((n_k, n_z)); Vn = np.empty_like(V)
while True:
    EV = V @ Pi.T          # (n_k,n_z): one matmul
    for j in range(n_z):
        g[:, j]  = golden_section(j, EV)
        Vn[:, j] = rhs(g[:, j], y[:, j], EV[:, j])
    d = np.max(np.abs(Vn - V))
    V[:] = Vn              # copy, do not rebind

    for _ in range(20):    # Howard: no argmax
        EV = V @ Pi.T
        for j in range(n_z):
            V[:,j] = rhs(g[:,j], y[:,j], EV[:,j])
    if d < 1e-6:
        break
\end{lstlisting}
\end{column}
\begin{column}{0.44\textwidth}
\small
\texttt{V @ Pi.T} is the whole expectation --- $\mathbb{E}[v(k',z')\mid z]$ for every $(k',z)$ at once. It is computed \textbf{once per sweep}, outside the maximization, not re-summed inside it.
\medskip
\begin{conceptbox}\footnotesize
The inner block is Howard: the policy $g$ is \emph{held fixed} and only the value is updated, so there is no \texttt{argmax} and no bracket --- just an interpolation and an add.
\end{conceptbox}
\medskip
\begin{examplebox}\footnotesize
Two nested loops, and they do completely different work. Confusing them is the most common way this algorithm is implemented wrongly.
\end{examplebox}
\end{column}
\end{columns}
\end{frame}

%=============================================================================
\begin{frame}{Step 5: What It Bought, Stage by Stage}
\begin{columns}[T]
\begin{column}{0.55\textwidth}
\small
Same model, same grid ($n_k=300$), same chain, same tolerance --- only the inner maximization and the sweep structure differ.
\medskip
\begin{center}\footnotesize
\begin{tabularx}{\linewidth}{@{}>{\raggedright\arraybackslash}X r r@{}}
\toprule
 & \textbf{3.A4 style} & \textbf{this deck} \\
\midrule
policy updates & 1{,}352 & \textbf{66} \\
wall time & 1.03\,s & \textbf{0.93\,s} \\
inner work & 852M comparisons & 11.4M evals \\
worst $\log_{10}\mathcal{E}$ & $-1.30$ & $\mathbf{-2.52}$ \\
mean $\log_{10}\mathcal{E}$ & $-2.14$ & $\mathbf{-3.40}$ \\
distinct $k'$ / 36 states & 35 & \textbf{36} \\
\bottomrule
\end{tabularx}
\end{center}
\medskip
\footnotesize
\textbf{1.2 decades of accuracy, a twentieth of the policy updates, and inner work down by a factor of 75} --- at slightly \emph{less} wall clock. \src{4.A2 and 4.B4 report the same experiment at $n_k=500$ ($-1.55$ and $-2.77$); the errors there are smaller because the grid is finer. The gain is 1.2 decades either way.}
\end{column}
\begin{column}{0.42\textwidth}
\begin{takeawaybox}[title=\textbf{Where each gain came from}]\small
The \textbf{iteration count} is Howard. The \textbf{accuracy} is golden section plus interpolation. The \textbf{chain} changed nothing about speed and everything about what the model is.
\end{takeawaybox}
\medskip
\begin{cautionbox}\footnotesize
\textbf{But do not read too much into the clock.} The 75$\times$ drop in evaluations buys only a 10\% drop in time, because golden section runs in a Python loop while the grid search is one array operation. On 4.B4's finer grid the same comparison runs \emph{slower}. Report the operation count, which is implementation-independent, \emph{and} the clock.
\end{cautionbox}
\end{column}
\end{columns}
\end{frame}

%=============================================================================
\begin{frame}{The Staircase, Retired}
\begin{columns}[c]
\begin{column}{0.52\textwidth}
\centering
\begin{tikzpicture}
\begin{axis}[
    width=7.6cm, height=5.6cm, xmin=25.4, xmax=29.0, ymin=-0.035, ymax=0.115,
    xlabel={\scriptsize capital $k$ (36 consecutive grid points)},
    ylabel={\scriptsize saving $g(k)-k$},
    tick label style={font=\scriptsize}, scaled ticks=false,
    x tick label style={/pgf/number format/fixed, /pgf/number format/precision=1},
    y tick label style={/pgf/number format/fixed, /pgf/number format/precision=2},
    axis x line*=bottom, axis y line*=left,
    legend style={font=\tiny, draw=none, fill=none, at={(0.98,0.97)}, anchor=north east}]
  % coordinates are (k, k') from the ledger; pgfplots forms k'-k
  \addplot[thick, black!45, mark=square*, mark size=1.1pt]
    table[row sep=\\, x index=0, y expr=\thisrowno{1}-\thisrowno{0}] {
      25.55150 25.64631 \\
      25.64631 25.74112 \\
      25.74112 25.83593 \\
      25.83593 25.93074 \\
      25.93074 26.02556 \\
      26.02556 26.12037 \\
      26.12037 26.21518 \\
      26.21518 26.30999 \\
      26.30999 26.40480 \\
      26.40480 26.49961 \\
      26.49961 26.59442 \\
      26.59442 26.68923 \\
      26.68923 26.78404 \\
      26.78404 26.87885 \\
      26.87885 26.97366 \\
      26.97366 27.06847 \\
      27.06847 27.16328 \\
      27.16328 27.25810 \\
      27.25810 27.35291 \\
      27.35291 27.35291 \\
      27.44772 27.44772 \\
      27.54253 27.54253 \\
      27.63734 27.63734 \\
      27.73215 27.73215 \\
      27.82696 27.82696 \\
      27.92177 27.92177 \\
      28.01658 28.01658 \\
      28.11139 28.11139 \\
      28.20620 28.20620 \\
      28.30101 28.30101 \\
      28.39582 28.39582 \\
      28.49064 28.49064 \\
      28.58545 28.58545 \\
      28.68026 28.68026 \\
      28.77507 28.77507 \\
      28.86988 28.86988 \\
    };
  \addlegendentry{grid search (3.A4)}
  \addplot[very thick, RedTitle, mark=*, mark size=0.9pt]
    table[row sep=\\, x index=0, y expr=\thisrowno{1}-\thisrowno{0}] {
      25.55150 25.65455 \\
      25.64631 25.74556 \\
      25.74112 25.83660 \\
      25.83593 25.93074 \\
      25.93074 26.02370 \\
      26.02556 26.11484 \\
      26.12037 26.20596 \\
      26.21518 26.29708 \\
      26.30999 26.38820 \\
      26.40480 26.47932 \\
      26.49961 26.57043 \\
      26.59442 26.66153 \\
      26.68923 26.75262 \\
      26.78404 26.84371 \\
      26.87885 26.93480 \\
      26.97366 27.02588 \\
      27.06847 27.11695 \\
      27.16328 27.20799 \\
      27.25810 27.29899 \\
      27.35291 27.39008 \\
      27.44772 27.48119 \\
      27.54253 27.57230 \\
      27.63734 27.66341 \\
      27.73215 27.75452 \\
      27.82696 27.84562 \\
      27.92177 27.93671 \\
      28.01658 28.02780 \\
      28.11139 28.11888 \\
      28.20620 28.20996 \\
      28.30101 28.30101 \\
      28.39582 28.39582 \\
      28.49064 28.48693 \\
      28.58545 28.57799 \\
      28.68026 28.66906 \\
      28.77507 28.76013 \\
      28.86988 28.85120 \\
    };
  \addlegendentry{golden section}
  \draw[black!55, dashed] (axis cs:25.4,0) -- (axis cs:29.0,0);
  \node[font=\tiny, text=black!60, anchor=north west] at (axis cs:25.5,-0.004) {$g(k)=k$};
\end{axis}
\end{tikzpicture}
\end{column}
\begin{column}{0.45\textwidth}
\small
Thirty-six consecutive grid points in the middle state, plotted as \emph{saving} rather than $k'$ --- the same data, with the 45-degree line subtracted off.
\medskip
\begin{itemize}\footnotesize
  \item Grid search can only pick $k'=k+mh$, so saving takes \textbf{two values in this whole window}: one grid spacing, then zero. It is a staircase with two steps.
  \item Golden section gives a \textbf{smooth, declining} line that crosses zero once.
\end{itemize}
\medskip
\begin{cautionbox}\footnotesize
And the two cross zero a full \textbf{ten grid spacings apart} --- 27.35 against $k^{*}=28.35$. \textbf{4.A2} had to make this point with the residual, because the plot of $k'$ against $k$ hides it entirely. Subtract the 45-degree line and it is unmissable.
\end{cautionbox}
\end{column}
\end{columns}
\end{frame}

%=============================================================================
\begin{frame}{The Policy, and the Three Checks}
\begin{columns}[c]
\begin{column}{0.52\textwidth}
\centering
\begin{tikzpicture}
\begin{axis}[
    width=7.6cm, height=5.6cm, xmin=14.0, xmax=42.3, ymin=-0.75, ymax=0.68,
    xlabel={\scriptsize capital $k$}, ylabel={\scriptsize saving $g(k,z)-k$},
    tick label style={font=\scriptsize}, scaled ticks=false,
    y tick label style={/pgf/number format/fixed, /pgf/number format/precision=2},
    axis x line*=bottom, axis y line*=left,
    legend style={font=\tiny, draw=none, fill=none, at={(0.02,0.04)}, anchor=south west}]
  % (k, k') from the ledger; pgfplots forms k'-k, i.e. the 45-degree line subtracted off
  \addplot[very thick, RedTitle]
    table[row sep=\\, x index=0, y expr=\thisrowno{1}-\thisrowno{0}] {
      14.1742 14.7692 \\
      14.7431 15.3211 \\
      15.3119 15.8786 \\
      15.8808 16.4299 \\
      16.4497 16.9808 \\
      17.0185 17.5315 \\
      17.5874 18.0818 \\
      18.1563 18.6318 \\
      18.7251 19.1872 \\
      19.2940 19.7368 \\
      19.8629 20.2862 \\
      20.4317 20.8353 \\
      21.0006 21.3842 \\
      21.5694 21.9383 \\
      22.1383 22.4869 \\
      22.7072 23.0353 \\
      23.2760 23.5836 \\
      23.8449 24.1316 \\
      24.4138 24.6847 \\
      24.9826 25.2326 \\
      25.5515 25.7803 \\
      26.1204 26.3279 \\
      26.6892 26.8789 \\
      27.2581 27.4275 \\
      27.8270 27.9748 \\
      28.3958 28.5221 \\
      28.9647 29.0690 \\
      29.5336 29.6205 \\
      30.1024 30.1676 \\
      30.6713 30.7141 \\
      31.2401 31.2611 \\
      31.8090 31.8090 \\
      32.3779 32.3585 \\
      32.9467 32.9052 \\
      33.5156 33.4516 \\
      34.0845 33.9980 \\
      34.6533 34.5486 \\
      35.2222 35.0947 \\
      35.7911 35.6409 \\
      36.3599 36.1870 \\
      36.9288 36.7375 \\
      37.4977 37.2833 \\
      38.0665 37.8292 \\
      38.6354 38.3751 \\
      39.2043 38.9209 \\
      39.7731 39.4709 \\
      40.3420 40.0165 \\
      40.9108 40.5622 \\
      41.4797 41.1079 \\
      42.0486 41.6577 \\
    };
  \addlegendentry{$z$ high}
  \addplot[very thick, orange!90!black]
    table[row sep=\\, x index=0, y expr=\thisrowno{1}-\thisrowno{0}] {
      14.1742 14.6669 \\
      14.7431 15.2174 \\
      15.3119 15.7735 \\
      15.8808 16.3234 \\
      16.4497 16.8730 \\
      17.0185 17.4223 \\
      17.5874 17.9713 \\
      18.1563 18.5257 \\
      18.7251 19.0744 \\
      19.2940 19.6228 \\
      19.8629 20.1710 \\
      20.4317 20.7189 \\
      21.0006 21.2720 \\
      21.5694 21.8198 \\
      22.1383 22.3673 \\
      22.7072 22.9147 \\
      23.2760 23.4657 \\
      23.8449 24.0140 \\
      24.4138 24.5610 \\
      24.9826 25.1080 \\
      25.5515 25.6545 \\
      26.1204 26.2060 \\
      26.6892 26.7526 \\
      27.2581 27.2990 \\
      27.8270 27.8456 \\
      28.3958 28.3958 \\
      28.9647 28.9423 \\
      29.5336 29.4886 \\
      30.1024 30.0346 \\
      30.6713 30.5807 \\
      31.2401 31.1309 \\
      31.8090 31.6766 \\
      32.3779 32.2224 \\
      32.9467 32.7681 \\
      33.5156 33.3182 \\
      34.0845 33.8636 \\
      34.6533 34.4091 \\
      35.2222 34.9545 \\
      35.7911 35.5043 \\
      36.3599 36.0495 \\
      36.9288 36.5948 \\
      37.4977 37.1400 \\
      38.0665 37.6873 \\
      38.6354 38.2346 \\
      39.2043 38.7796 \\
      39.7731 39.3246 \\
      40.3420 39.8697 \\
      40.9108 40.4189 \\
      41.4797 40.9637 \\
      42.0486 41.5086 \\
    };
  \addlegendentry{$z$ middle}
  \addplot[very thick, black!55]
    table[row sep=\\, x index=0, y expr=\thisrowno{1}-\thisrowno{0}] {
      14.1742 14.5710 \\
      14.7431 15.1223 \\
      15.3119 15.6749 \\
      15.8808 16.2236 \\
      16.4497 16.7719 \\
      17.0185 17.3199 \\
      17.5874 17.8718 \\
      18.1563 18.4210 \\
      18.7251 18.9686 \\
      19.2940 19.5159 \\
      19.8629 20.0628 \\
      20.4317 20.6150 \\
      21.0006 21.1618 \\
      21.5694 21.7086 \\
      22.1383 22.2552 \\
      22.7072 22.8020 \\
      23.2760 23.3526 \\
      23.8449 23.8987 \\
      24.4138 24.4450 \\
      24.9826 24.9912 \\
      25.5515 25.5413 \\
      26.1204 26.0873 \\
      26.6892 26.6327 \\
      27.2581 27.1786 \\
      27.8270 27.7288 \\
      28.3958 28.2739 \\
      28.9647 28.8193 \\
      29.5336 29.3646 \\
      30.1024 29.9128 \\
      30.6713 30.4594 \\
      31.2401 31.0045 \\
      31.8090 31.5495 \\
      32.3779 32.0945 \\
      32.9467 32.6438 \\
      33.5156 33.1886 \\
      34.0845 33.7334 \\
      34.6533 34.2782 \\
      35.2222 34.8273 \\
      35.7911 35.3718 \\
      36.3599 35.9164 \\
      36.9288 36.4610 \\
      37.4977 37.0099 \\
      38.0665 37.5542 \\
      38.6354 38.0986 \\
      39.2043 38.6431 \\
      39.7731 39.1917 \\
      40.3420 39.7360 \\
      40.9108 40.2802 \\
      41.4797 40.8245 \\
      42.0486 41.3729 \\
    };
  \addlegendentry{$z$ low}
  \draw[black!45, dashed] (axis cs:14.0,0) -- (axis cs:42.3,0);
  \draw[black!45, dashed] (axis cs:28.3484,-0.75) -- (axis cs:28.3484,0.68);
  \node[font=\tiny, text=black!60, anchor=south west, rotate=90] at (axis cs:28.3484,-0.70) {$k^{*}=28.35$};
\end{axis}
\end{tikzpicture}
\end{column}
\begin{column}{0.45\textwidth}
\small
\textbf{Three things to look at, in this order:}
\medskip
\begin{enumerate}\footnotesize
  \item The three curves are \textbf{ordered in $z$} and do not touch --- a good shock saves more at every $k$. A transposed $\Pi$ breaks this immediately.
  \item Saving is \textbf{declining} in $k$ and nearly linear, which is the concavity the theory promises.
  \item The middle curve crosses zero at $k^{*}=28.35$; the low and high curves cross on either side. Those are the model's conditional steady states.
\end{enumerate}
\medskip
\begin{cautionbox}\footnotesize
Note how \emph{shallow} these curves are --- saving moves by about $1$ unit while $k$ moves by $28$. That is why the previous slide's grid-search crossing missed by ten spacings: a flat function makes a weak test. \textbf{Use the residual to measure accuracy}; use the picture for sign checks only.
\end{cautionbox}
\end{column}
\end{columns}
\end{frame}

%=============================================================================
\begin{frame}{What Is Still Crude}
\begin{columns}[T]
\begin{column}{0.53\textwidth}
\small
Even with the whole Session 4 toolkit switched on, three things in this code are still the cheap option:
\medskip
\begin{itemize}\footnotesize
  \item \textbf{Linear interpolation.} $O(h^{2})$, and the kinked derivative is felt by anything that differentiates $\hat v$. A shape-preserving spline would do better without losing concavity.
  \item \textbf{One continuous state.} Add labour or a second asset and the grid is $n_k^{2}$ --- and \textbf{4.B1's Smolyak} is the first answer, S5's projection the second.
  \item \textbf{We still iterate on the value function.} \textbf{4.B4's EGM} skips it entirely: $10^{-6.06}$ against this solver's $10^{-2.77}$ on the same grid --- \textbf{three further decades}, and 38 times faster.
\end{itemize}
\end{column}
\begin{column}{0.44\textwidth}
\begin{takeawaybox}[title=\textbf{The honest summary}]\small
This is a good solver. It is not the best one in the session --- EGM is --- and the reason to write this one first is that \textbf{every piece of it generalises} to models where EGM does not apply.
\end{takeawaybox}
\medskip
\begin{examplebox}\footnotesize
EGM needs an invertible $u'$, one continuous control and a monotone policy. Value iteration with golden section needs none of those, which is why it is still the workhorse for models with discrete choices or binding constraints.
\end{examplebox}
\end{column}
\end{columns}
\end{frame}

%=============================================================================
\begin{frame}{Takeaways}
\begin{takeawaybox}
\begin{itemize}\small
  \item The same five steps as 3.A4, with four substitutions: \textbf{Rouwenhorst, interpolation, golden section, Howard}. The code is barely longer.
  \item \textbf{Check the chain first.} Implied $\rho=0.950000$ and $\sigma_z$ matching the target to eight digits --- three lines that rule out the commonest silent error in the literature.
  \item The interpolation choice is a \textbf{hypothesis of the maximizer}: linear keeps $\hat v$ concave, concavity gives unimodality, and unimodality is what licenses golden section. A cubic spline breaks the chain at the first link.
  \item Measured: \textbf{66 policy updates against 1,352}, inner work down $75\times$, and \textbf{1.2 decades} of Euler accuracy. The wall clock barely moves, and its sign flips with the grid size --- which is why you report the operation count too.
  \item It is still not the best solver here. \textbf{EGM is}, by three further decades --- but this one survives into models with discrete choices and binding constraints, where EGM cannot go.
\end{itemize}
\end{takeawaybox}
\end{frame}

%=============================================================================
\begin{frame}{Bridge: Two Sessions, One Model, Many Solvers}
\begin{columns}[T]
\begin{column}{0.52\textwidth}
\small
Sessions 3 and 4 have now solved the same model many different ways --- grid VFI, time iteration, Howard, monotonicity, continuous choice, EGM --- and every difference between them was measured rather than asserted. That is the habit the rest of the course rests on.
\medskip

What has not changed is the \textbf{representation}: values, on a grid, one number per state. Everything we have built makes that representation cheaper or more accurate. None of it makes the grid go away --- and at two continuous states it already hurts.
\end{column}
\begin{column}{0.45\textwidth}
\begin{conceptbox}[title=\textbf{Next week (S5)}]\small
\textbf{Perturbation} discards the grid and expands around the steady state. \textbf{Projection} replaces $N$ separate values with a handful of basis coefficients --- and Smolyak returns, now as the basis rather than the node set.\\[1mm]
Then Aiyagari, end to end.
\end{conceptbox}
\medskip
\begin{examplebox}\footnotesize
\textbf{Homework 1} is due October 22. Part B Q3 is this deck's model, and asks you to build the chain both ways and see whether it moved the answer.
\end{examplebox}
\end{column}
\end{columns}
\end{frame}

\end{document}
